

\section{Summary of the Project}

The Medicare Bot project successfully demonstrates the integration of artificial intelligence, embedded systems, and web technologies to create a smart medication dispensing system. The system is capable of identifying patients using face recognition, verifying medication schedules, and automatically dispensing the correct medication.
The system design, implementation, and evaluation discussed in Chapters \ref{systemarchitecture} to \ref{result} collectively demonstrate the effectiveness of the proposed solution.
The architecture combines a Raspberry Pi-based AI processing unit, an Arduino-based motor control system, and a full-stack software solution consisting of a FastAPI backend and a React frontend. This integration enables real-time monitoring, data logging, and visualization of dispensing activities.

\section{Key Achievements}

The major achievements of the project include:

\begin{itemize}
    \item Implementation of a face recognition system using deep learning models
    \item Development of an automated medication dispensing mechanism
    \item Integration of Raspberry Pi and Arduino for coordinated hardware control
    \item Design and implementation of a backend API for data management
    \item Development of a frontend dashboard for monitoring system activity
    \item Implementation of image-based logging for verification of dispensing events
\end{itemize}

\section{System Effectiveness}

The system operates as a functional prototype capable of performing end-to-end operations, including patient recognition, medication dispensing, and real-time logging. The use of automation reduces the possibility of human error and improves medication adherence.

The system demonstrates reliable performance within the constraints of the Raspberry Pi hardware, and the implemented optimizations ensure efficient operation.

\section{Practical Significance}

The Medicare Bot has potential applications in hospitals, elderly care centers, and home healthcare environments. It can assist healthcare providers by automating routine tasks, reducing workload, and improving the accuracy of medication delivery.

The modular design also allows future expansion and adaptation for different use cases such as delivery of medical supplies or integration with hospital management systems.

\section{Future Scope}

Although the system is functional, several improvements can be made in future work:

\begin{itemize}
    \item Implementation of a mobile or web-based control panel for remote management
    \item Enhancement of face recognition accuracy using more advanced models
    \item Addition of voice interaction for improved user experience
    \item Deployment on cloud infrastructure for scalability
    \item Implementation of missed dose detection and notification system
\end{itemize}

Integration with cloud-based healthcare systems and IoT platforms can further enhance scalability and enable real-time remote patient monitoring.

\section{Final Remarks}

In conclusion, the Medicare Bot project highlights the potential of combining artificial intelligence and embedded systems to solve real-world healthcare challenges. The system provides a foundation for future smart healthcare solutions and demonstrates how automation can enhance patient care, safety, and operational efficiency.